\documentclass[11pt,a4paper]{article}
\usepackage{kotex}
\usepackage{fontspec}
\setmainfont{Times New Roman}
\setmainhangulfont{AppleMyungjo}
\setsanshangulfont{Apple SD Gothic Neo}
\usepackage[a4paper,top=18mm,bottom=18mm,left=20mm,right=20mm]{geometry}
\usepackage{fancyhdr}
\setlength{\headheight}{24pt}
\usepackage{enumitem}
\usepackage{mdframed}
\usepackage{booktabs}
\linespread{1.18}
\setlength{\parindent}{0pt}
\pagestyle{fancy}
\fancyhf{}
\renewcommand{\headrulewidth}{0.6pt}
\fancyhead[L]{\sffamily\bfseries 2026 고1 영어 변형 - 정답 및 해설 지문 43-45}
\fancyhead[R]{\sffamily Emma와 Lisa}
\fancyfoot[C]{\framebox[16mm][c]{\thepage}}
\newcommand{\qno}[1]{\par\vspace{10pt}{\sffamily\bfseries #1.}\ }
\newmdenv[linewidth=0.6pt,roundcorner=3pt,innertopmargin=7pt,innerbottommargin=7pt,
          innerleftmargin=9pt,innerrightmargin=9pt,skipabove=5pt,skipbelow=5pt]{pbox}
\begin{document}
\begin{center}
{\fontsize{15}{17}\selectfont\sffamily\bfseries [지문 43-45] 정답 및 해설}
\end{center}
\vspace{6pt}
\begin{center}
\renewcommand{\arraystretch}{1.35}
\begin{tabular}{cccc}
\toprule
문항 & 유형 & 정답 & 배점 \\
\midrule
1 & 글 순서 배열 & ② & 2점 \\
2 & 지칭 추론 & ② & 2점 \\
3 & 내용 불일치 & ④ & 2점 \\
4 & 서술형 심경 변화 & 모범답안 참조 & 2점 \\
\bottomrule
\end{tabular}
\end{center}

\qno{1}\textbf{글 순서 배열 - 정답: ② (C)-(B)-(D)}
\begin{pbox}
(A)에서 Emma가 대회에 나가기로 결심한다. 이어 (C)에서 실제 대회 당일의 상황, Lisa와의 비교, Emma의 공연, 그리고 Emma의 1등 결과가 제시된다. 그 다음 (B)에서 Lisa가 Emma에게 다가와 축하와 격려의 말을 건넨다. 마지막으로 (D)의 ``These words''는 바로 (B)의 Lisa의 말을 가리키므로 (D)는 (B) 뒤에 와야 한다.
\end{pbox}

\qno{2}\textbf{지칭 추론 - 정답: ②}
\begin{pbox}
(a), (c), (d), (e)는 모두 Emma를 가리킨다. 반면 (b) \textit{She didn't seem bothered by the results}는 Emma를 축하하러 온 Lisa가 결과에 신경 쓰지 않는 모습을 묘사하므로 Lisa를 가리킨다. 따라서 나머지 넷과 다른 것은 ②이다.
\end{pbox}

\qno{3}\textbf{내용 불일치 - 정답: ④}
\begin{pbox}
본문은 Emma가 ``first place''를 차지했다고 명시한다. 따라서 ④ ``runner-up position''은 2등을 의미하므로 본문과 일치하지 않는다.

\medskip
①은 같은 고등학교 출신이라는 내용, ②는 유명해진 Lisa가 Emma에게 따뜻한 말을 건넨 내용, ③은 Emma가 공연장에 들어가 Lisa를 응원하는 관중을 본 내용, ⑤는 이후 두 사람이 같은 무대에서 노래했다는 내용과 일치한다.
\end{pbox}

\qno{4}\textbf{서술형 심경 변화}
\begin{pbox}
\textbf{모범답안}: Emma는 대회 전 Lisa의 성공을 보며 질투와 긴장을 느꼈지만, 대회에서 1등을 하고 Lisa의 진심 어린 축하를 받은 뒤 부끄러움과 감사함을 느끼며 더 나은 방향으로 나아갈 용기를 얻었다.

\medskip
\textbf{채점 기준}: 대회 전 심경(질투, 시기, 긴장, 비교 의식 등)과 대회 후 심경(부끄러움, 감사, 감동, 자신감, 용기 등)을 각각 하나 이상 포함하고 변화의 방향을 한 문장으로 설명하면 정답.
\end{pbox}
\end{document}
